
\emph{suppose que pour tout point \(s\in S\), et toute composante irr\'eductible \(Y_s^i\) de \(Y_s\), les composantes irr\'eductibles de \(X_s\cap Y_s^i=V(\varphi|Y_s^i)\) soient de dimension \(\ge2\), (ce qui est le cas si les composantes irr\'eductibles des fibres de \(Y\) sont de dimension \(\ge3\)). Alors le morphisme induit entre les foncteurs de Picard}
\[
f^*:\underline{\Pic}_{Y/S}\longrightarrow\underline{\Pic}_{X/S}
\]
\emph{est de type fini.}

\emph{De plus, si \(S\) est int\`egre, il existe un ouvert non-vide \(U\) de \(S\) tel que les foncteurs de Picard restreints \`a \(U\) soient repr\'esentables et le morphisme induit}
\[
\underline{\Pic}_{Y/S}|U\longrightarrow\underline{\Pic}_{X/S}|U
\]
\emph{soit quasi-affine (de type fini).}

D\'emontrons d'abord la premi\`ere assertion. Comme dans la preuve de 3.6, on voit par r\'ecurrence noeth\'erienne qu'on peut remplacer \(S\) par un ouvert non-vide quelconque, et on se ram\`ene facilement au cas o\`u \(S\) est int\`egre et \(Y\) est int\`egre sur \(S\), toujours en rempla\c{c}ant \(X\) par \(X\times_YY'\) lorsqu'on remplace \(Y\) par \(Y'\)\,; ceci fait, on se ram\`ene de la m\^eme fa\c{c}on au cas o\`u \(Y\) est normal sur \(S\), en utilisant \`a la place de XII 1.3 le lemme suivant.

\textbf{Lemme 3.9.} \emph{Soient \(S\) un sch\'ema noeth\'erien int\`egre et \(X\) un \(S\)-sch\'ema, int\`egre sur \(S\). Alors il existe un ouvert non-vide \(U\) de \(S\), un sch\'ema int\`egre \(S'\) fini et fid\`element plat sur \(U\), un \(S'\)-sch\'ema \(X'\) normal et int\`egre sur \(S'\), et un \(S'\)-morphisme fini surjectif \(X'\to X_{S'}\).}

